\section{The fixed-frozen principal obstruction}
\label{sec:principal-obstruction}

The mutable globalization theorem is valid in every rank.  A full principal
quantum seed contains an additional frozen lattice, and requiring one fixed
frozen basis on the Hall side is strictly stronger.  This section gives the
complete criterion.

\subsection{The enlarged term--frozen lattice}

Let $F\simeq\mathbb Z^n$ be the principal frozen lattice and put
\[
 \widehat L=L\oplus F.
\]
The frozen basis is fixed once and for all.  For a maximal rigid chart $R$,
define the lift
\begin{equation}
 J_R=
 \begin{pmatrix}
  \Xi_R&0\\
  0&I_n
 \end{pmatrix}:
 K\oplus F\longrightarrow L\oplus F.
 \label{eq:J-R}
\end{equation}
The first block sends the mutable seed basis to the reduced exact term classes;
the second identifies the frozen basis in every chart.

\begin{definition}[Fixed-frozen globalization]
\label{def:full-globalization}
Let $S\in\Skew_n(\mathbb Z)$ and $d>0$.  A
\emph{contractible-radical fixed-frozen globalization} of the affine
principal quantisation \cref{eq:principal-affine-family} is an integral skew
form $\widehat{\widetilde\Lambda}_S$ on $\widehat L$ such that:
\begin{enumerate}[label=(\roman*),leftmargin=2.3em]
\item $C\oplus0$ lies in its radical;
\item its pullback along $J_R$ is
      $\widetilde\Lambda_R(S)$ for every reachable maximal rigid chart $R$.
\end{enumerate}
\end{definition}

The condition at the initial chart already determines the form on the
quotient
\[
 \widehat L/(C\oplus0)\simeq K\oplus F.
\]
To display it, set
\begin{equation}
 C_S:=dI_n+SB_0,
 \qquad
 R_S:=dB_0^T+B_0^TSB_0.
 \label{eq:C-S-R-S}
\end{equation}

\begin{lemma}[Initial quotient form]
\label{lem:initial-quotient-full}
If a fixed-frozen globalization exists, then its quotient form on $K\oplus F$
is uniquely
\begin{equation}
 \overline\Lambda_S=
 \begin{pmatrix}
  S&C_S\\
  -C_S^T&R_S
 \end{pmatrix}
 =
 \begin{pmatrix}
  S&dI+SB_0\\
  -dI+B_0^TS&dB_0^T+B_0^TSB_0
 \end{pmatrix}.
 \label{eq:quotient-full-form}
\end{equation}
The candidate global form is the pullback of
$\overline\Lambda_S$ along
\[
 \widehat D=\begin{pmatrix}D&0\\0&I_n\end{pmatrix}:
 L\oplus F\longrightarrow K\oplus F.
\]
\end{lemma}

\begin{proof}
At the initial morphism chart,
$\Delta_0=-I_n$, $G_0=I_n$, and the reduced term lift is the $Z$-boundary.
Expanding \cref{eq:principal-affine-family} gives
\[
 \widetilde\Lambda_0(S)=
 \begin{pmatrix}
  S&-dI-SB_0\\
  dI-B_0^TS&dB_0^T+B_0^TSB_0
 \end{pmatrix}.
\]
The quotient mutable coordinate is the split index, equal to minus the
$Z$-boundary coordinate.  Conjugating by
$\operatorname{diag}(-I,I)$ gives exactly
\cref{eq:quotient-full-form}.  The radical condition forces factorization
through $\widehat D$, proving uniqueness.
\end{proof}

\subsection{The complete criterion}

\begin{theorem}[Fixed-frozen principal globalization criterion]
\label{thm:full-globalization-criterion}
Let $S\in\Skew_n(\mathbb Z)$.  A contractible-radical fixed-frozen
globalization exists if and only if, for every reachable maximal rigid chart
$R$, both
\begin{align}
 \Delta_R^TS\Delta_R&=G_R^TSG_R,
 \label{eq:full-mutable-condition}\\
 (\Delta_R^T+G_R^T)(dI_n+SB_0)&=0
 \label{eq:full-mixed-condition}
\end{align}
hold.  Equivalently, $S\in\Qglob$ and the mixed-block equations
\cref{eq:full-mixed-condition} hold.  When it exists, the global form is
unique and equals
\begin{equation}
 \widehat{\widetilde\Lambda}_S
 =\widehat D^T\overline\Lambda_S\widehat D.
 \label{eq:full-global-form}
\end{equation}
\end{theorem}

\begin{proof}
By \cref{lem:initial-quotient-full}, only the form
\cref{eq:full-global-form} can occur.  Pulling it back along $J_R$ gives
\begin{equation}
 J_R^T\widehat{\widetilde\Lambda}_SJ_R
 =\begin{pmatrix}
   \Delta_R^TS\Delta_R&\Delta_R^TC_S\\
   -C_S^T\Delta_R&R_S
  \end{pmatrix}.
 \label{eq:actual-full-pullback}
\end{equation}
On the other hand, expanding \cref{eq:principal-affine-family} at $R$ gives
\begin{equation}
 \widetilde\Lambda_R(S)
 =\begin{pmatrix}
   G_R^TSG_R&-G_R^TC_S\\
   C_S^TG_R&R_S
  \end{pmatrix}.
 \label{eq:desired-full-form}
\end{equation}
Equality of the upper-left blocks is
\cref{eq:full-mutable-condition}.  Equality of the upper-right blocks is
\[
 \Delta_R^TC_S=-G_R^TC_S,
\]
which is \cref{eq:full-mixed-condition}.  The lower-left equation is its
negative transpose, and the lower-right blocks already agree.  This proves
necessity and sufficiency.
\end{proof}

\begin{definition}[Fixed-frozen parameter locus]
The full fixed-frozen parameter locus is the affine integral set
\begin{equation}
 \mathfrak Q_T^{\mathrm{ff}}
 :=\left\{
 S\in\Qglob:
 (\Delta_R^T+G_R^T)(dI+SB_0)=0
 \text{ for every }R
 \right\}.
 \label{eq:Q-ff}
\end{equation}
It is cut out by an inhomogeneous integer linear system on the entries of
$S$.
\end{definition}

\subsection{Canonical principal no-go and full-rank recovery}

\begin{corollary}[Invertible mixed block]
\label{cor:invertible-mixed-block}
If $dI+SB_0$ is invertible over $\mathbb Q$, then a fixed-frozen
globalization can exist only if
\begin{equation}
 \Delta_R=-G_R
 \label{eq:Delta-minus-G}
\end{equation}
for every reachable chart $R$.
\end{corollary}

\begin{proof}
Right multiplication of \cref{eq:full-mixed-condition} by
$(dI+SB_0)^{-1}$ gives
$\Delta_R^T+G_R^T=0$.
\end{proof}

\begin{corollary}[Canonical principal fixed-frozen obstruction]
\label{cor:canonical-full-no-go}
For the canonical principal parameter $S=0$, the mutable Hall cocycle always
globalizes by \cref{cor:zero-globalizes}.  The full fixed-frozen principal
form globalizes if and only if
\[
 \Delta_R=-G_R
\]
for every reachable chart.  In particular, one chart with
$\Delta_R\ne-G_R$ is a complete obstruction.
\end{corollary}

\begin{proof}
For $S=0$, the matrix $dI+SB_0=dI$ is invertible.  Apply
\cref{thm:full-globalization-criterion,cor:invertible-mixed-block}.
\end{proof}

The canonical principal coefficients therefore remove the rank obstruction to
constructing quantum seeds, and they remove the rank obstruction to a global
\emph{mutable} Hall cocycle, but they do not automatically produce one fixed
frozen Hall lattice.

The full-rank compatible parameter behaves differently.

\begin{corollary}[Full-rank fixed-frozen globalization]
\label{cor:full-rank-full-principal}
Assume $B_0^T\Lambda_0=dI_n$.  Then
\[
 \Lambda_0\in\mathfrak Q_T^{\mathrm{ff}}.
\]
More precisely,
\begin{equation}
 dI+\Lambda_0B_0=0,
 \qquad
 dB_0^T+B_0^T\Lambda_0B_0=0,
 \label{eq:full-rank-C-zero}
\end{equation}
and the quotient form \cref{eq:quotient-full-form} is
\begin{equation}
 \overline\Lambda_{\Lambda_0}
 =\begin{pmatrix}\Lambda_0&0\\0&0\end{pmatrix}.
 \label{eq:central-frozen-full-form}
\end{equation}
Thus the full principal fixed-frozen globalization is the coefficient-free
term form $D^T\Lambda_0D$ with a central frozen lattice.
\end{corollary}

\begin{proof}
Taking transposes in $B_0^T\Lambda_0=dI$ and using skew-symmetry gives
$\Lambda_0B_0=-dI$.  This proves the first equality in
\cref{eq:full-rank-C-zero}; the second follows from
\[
 B_0^T\Lambda_0B_0=dB_0.
\]
The mutable condition is \cref{eq:companion-identity}, while the mixed
condition is automatic because $C_{\Lambda_0}=0$.  Formula
\cref{eq:central-frozen-full-form} follows from
\cref{eq:quotient-full-form}.
\end{proof}

\begin{remark}[Two distinguished principal gauges]
The affine principal quantisation space contains two structurally important
parameters.
\begin{enumerate}[label=(\roman*),leftmargin=2.3em]
\item $S=0$ is canonical from the principal-coefficient viewpoint.  It always
      globalizes mutably but generally requires chamber-dependent frozen
      identifications.
\item In full rank, $S=\Lambda_0$ is the coefficient-free gauge.  It
      globalizes both mutably and with a fixed frozen lattice, but the frozen
      directions become central.
\end{enumerate}
The distinction is invisible if one records only the compatible extended
exchange matrices; it appears when one asks for a single Hall term lattice.
\end{remark}

\subsection{The crossed frozen atlas}

When $S\in\Qglob$ but $S\notin\mathfrak Q_T^{\mathrm{ff}}$, the mutable
chamber algebras are restrictions of the one global Hall twist
$\star_{\Omega_S}$, whereas the full principal forms must be transported with
chart-dependent frozen coordinates.  The transition from $R$ to $R'$ acts on
the quotient mutable lattice by
\[
 \mathsf L_{R'}\mathsf L_R^{-1}
\]
and on the frozen lattice by the principal coefficient transport encoded by
the $C$-matrices.  The resulting object is naturally a crossed atlas rather
than one fixed based quantum torus.

\begin{proposition}[Mutable/frozen separation]
\label{prop:mutable-frozen-separation}
For every $S\in\Qglob$:
\begin{enumerate}[label=(\roman*),leftmargin=2.3em]
\item one global exact-Hall cocycle controls all compatible rigid products;
\item the obstruction to one fixed principal frozen lattice is exactly the
      family of matrices
      \begin{equation}
       \mathcal O_R(S)
       :=(\Delta_R^T+G_R^T)(dI+SB_0);
       \label{eq:frozen-obstruction-matrix}
      \end{equation}
\item $S\in\mathfrak Q_T^{\mathrm{ff}}$ if and only if every
      $\mathcal O_R(S)$ vanishes.
\end{enumerate}
\end{proposition}

\begin{proof}
Part (i) is \cref{thm:globalizable-Hall-theta}; parts (ii) and (iii) are
\cref{thm:full-globalization-criterion}.
\end{proof}

\begin{algorithm}[Full principal globalization test]
\label{alg:full-principal-test}
For a finite rigid atlas and a chosen compatibility scalar $d$:
\begin{enumerate}[label=\arabic*.,leftmargin=2.2em]
\item introduce $\binom n2$ integer unknowns for a skew matrix $S$;
\item impose the homogeneous equations
      $\mathsf L_R^TS\mathsf L_R=S$;
\item impose the affine equations
      $(\Delta_R^T+G_R^T)(dI+SB_0)=0$;
\item compute the rational affine solution space and its integral points.
\end{enumerate}
The first subsystem computes $\Qglob$; the second cuts out
$\mathfrak Q_T^{\mathrm{ff}}$.
\end{algorithm}
